%! Author = Maximilian Fülle
%! Date = 29.05.2026


\chapter{Curves over valued fields and their models}\label{ch:curves-over-valued-fields-and-their-models}

\section{Basic Definitions}\label{sec:basic-definitions}
Let $K$ be a field endowed with an additive valuation of rank $1$, $v_K: K \to \hat{\RR} \coloneqq \RR\cup \lbrace\infty\rbrace$.
Let $\Gamma_K \coloneqq v_K(K^*) \subset \RR$ be the value group of $v_K$, $\OK \subset K$ its valuation ring with
maximal ideal $\mK$ and $k \coloneqq \OK /\mK$ its residue field.

For most of the first chapter, we make the following additional assumptions:
\begin{assumption} \label{ass:v_K}
    \begin{enumerate}[(a)]
        \item $v_K$ is discrete, i.e.\ $\Gamma_K$ is a discrete subgroup of $\RR$,
        \item $K$ is complete with respect to $v_K$,
        \item $k$ is perfect.
    \end{enumerate}
\end{assumption}

Throughout the whole monograph, $X$ will denote a \textit{nice} curve over $K$,
that is a smooth, projective and absolutely irreducible curve over $K$.
\textcolor{red}{If not stated otherwise, $\eta$ denotes the generic point of $X$.}

\begin{definition} \label{def:model}
    A \textit{model} $\X$ of $X$ is a proper, flat and normal $\OK $-scheme
    with an identification of its generic fibre $\X_g$ with $X$.
\end{definition}

\begin{example}
    Let $K = \QQ_p$, $\OK  = \ZZ_p$ and $X = \PP^1(K)$, then $\X = \PP^1(\OK )$ is a model of $X$ with the
    evident isomorphism from $\X_g = \X \times_{\OK} K$ to $\PP^1(K)$.
    The special fibre $\X_s$ has one irreducible component which is isomorphic to $\PP^1(k)$.
\end{example}

\begin{remark} \label{rem:flatIffIntegral}
    The condition of a model being flat is equivalent to it being integral:
    If $\X \to \Spec \OK $ is an integral scheme, then for every $x\in \X$, the stalk $\OO_{\X,x}$ is an integral ring extension of $\OK$.
    Thus $\OO_{\X,x}$ is a torsion-free module over the discrete valuation ring $\OK$, hence flat by~\cite[Corollary 1.2.14]{Liu}.
    Conversely, the integrality of a model $\X$ of $X$ follows from its flatness and the integrality of $X$ by~\cite[Proposition 4.3.8]{Liu}.
\end{remark}

\begin{definition}
    A \textit{modification} of a model of $X$ is a morphism $\X' \to \X$ of $\OK $-schemes which induces the
    identity on the generic fibre.
\end{definition}

\begin{remark}\label{rem:modifications-are-surjective-and-partial-order-of-models}
    \begin{enumerate}[(i)]
        \item Since a model is proper over $\OK $, a modification is a proper morphism.
            Thus since it is dominant, it is also surjective.
        \item We have a partial order of models of $X$ given by $\X \leq \X'$ if there exists a modification $\X' \to \X$.
    \end{enumerate}
\end{remark}

\begin{comment}
    \begin{example}
    In this example, we consider models of curves in the plane $\PP^2(K)$.
    Such a (smooth, projective and irreducible) curve $X$ is defined by a homogeneous irreducible polynomial $F(X,Y,Z) \in K[X,Y,Z]$.
    We show that if $F(X,Y,Z)$ is integral (that is all coefficients of $F$ are in $\OK$) and primitive, a model of $X$ is given by
    \begin{align}
        \X = \Proj{\OK[X,Y,Z]/(F)}.
    \end{align}
    Is is clearly a proper $\OK$-scheme.
    For the flatness, it suffices to  show that $R \coloneqq \OK[X,Y,Z]/(F(X,Y,Z))$ is a torsion-free $\OK$-module by Remark \ref{rem:flatIffIntegral}.
    This is the case if and only if $\bar{F} \neq 0 \in k[X,Y,Z]$, which is equivalent to $F$ being primitive. 
    For the normality, we use the following argument which is similar to Example 8.2.26 of \cite{Liu}.
    Since $\X$ is a Cohen-Macaulay scheme, it is normal if and only if it is normal at the points of codimension $1$.
    At this points, normality is equivalent to regularity.
    These are exactly the closed points on the generic fibre $\X_\eta=X$ and the generic points of the special fibre $\X_s$.
    Let $\xi \in \X$ be a generic point of the special fibre $\X_s$ and let $\pi\in \OK$ be a uniformising parameter.
    On the affine chart $U: Z \neq 0$ with $x = X/Z$ and $y=Y/Z$, the point $\xi$ corresponds to the ideal $\m$ that is generated by $\pi$ and a 
    polynomial $G(x,y) \in \OK[x,y]$ whose image $\bar{G}(x,y) \in k[x,y]$ is irreducible.
    We can write 
    \begin{align}
        F(x,y) = G(x,y)^r H_1(x,y) + \pi^e H_2(x,y),
    \end{align}
    for $H_i(x,y)\in \OK[x,y]$ and $r,e \geq 1$ and with $\bar{H}_1(x,y) \notin \bar{G}(x,y)k[x,y]$ and $\bar{H}_2(x,y) \neq 0$.
    By Corollary 4.2.12 of \cite{Liu}, $\xi$ is regular if $F \notin \m^2$.
    Thus, $\X$ is normal at $\xi$ if and only if either $r=1$; or $e=1$ and $\bar{H}_2(x,y) \neq 0$.
    \end{example}
\end{comment}

\begin{example}\label{ex:model-of-plane-curve}
    We discuss the important example of plane curves. 
    A smooth and absolutely irreducible curve in $\PP^2(K)$ is defined by a homogeneous irreducible polynomial $F \in K [X,Y,Z]$.
    Potentially after a normalisation, we can assume that $F$ is integral (i.e., all coefficients are in $\OK$) and primitive.
    Then the $\OK$-scheme 
    \begin{align}
        \X = \Proj{\OK[X,Y,Z]/(F)}
    \end{align}
    is an obvious candidate for a model of $X$.
    $\X$ is proper as a closed subscheme of $\PP^2(K)$ and flat by Remark \ref{rem:flatIffIntegral} since 
    $\OK[X,Y,Z]/(F)$ is torsion-free by irreducibility and primitivity of $F$.
    Hence, $\X$ is a model if and only if it is normal.
    The discussion of normality in this case is given in Example 8.2.26 of \cite{Liu} which we rephrase for our example:
    Since $\X$ is Cohen-Macualay, normality of $\X$ is equivalent to normality at the points of codimension $1$,
    which in turn is equivalent to regularity at these points.
    The points of codimension $1$ are the closed points of the (smooth) generic fibre $\X_\eta = X$,
    and the generic points of the special fibre $\X_s$.
    Now let $\xi \in \X_s$ be a generic point, then on the affine cover $Z \neq 0$
    with $x = X/Z$ and $y = Y/Z$, $\xi$ corresponds to the ideal $\m$ which is generated by a unifomising parameter $\pi \in \OK$
    and a polynomial $G \in K[x,y]$, whose reduction $\bar{G} \in k[x,y]$ is irreducible.
    We can write 
    \begin{align}
        F(x,y) = G(x,y)^r H_1(x,y) + \pi^e H_2(x,y),
    \end{align}
    where $\bar{H}_1(x,y) \notin \bar{G}(x,y)k[x,y]$ and $\bar{H}_2(x,y)\neq 0$.
    By \cite[Corollary 4.2.12]{Liu}, $\X$ is regular at $\xi$ if and only if $F \notin \m^2$,
    which in turn is the case if and only if either $r=1$; or $e=1$ and $\bar{H}_2 \notin \bar{G}(x,y)k[x,y]$.
    
\end{example}

\begin{example}\label{ex:particular-model-of-plane-curve}
    \begin{enumerate}
        \item Consider the curve $X$ defined by $y^2 = x^3 + px$ over $\QQ_p$ with $p \neq 2$.
            We have seen that $\X = \Proj{\OK[X,Y,Z]/(X^3 + pXY^2 - Y^2Z)}$ is a proper and flat $\OK$-scheme with $\X_\eta = X$.
            The special fibre is given by $\X_s = \Proj{k[X,Y,Z]/(X^3 - Y^2Z)}$.
            On the affine cover $Z\neq 0$, its generic point $\xi$ corresponds to the ideal $\m = (p, x^3-y^2)$.
            Thus, in the language of Example \ref{ex:model-of-plane-curve} it holds $r=e=1$, i.e. $\X$ is normal and therefore a model of $X$.
        \item Let $X: y^3 = 1 + x^3 + 3x$ be a curve over $\QQ_3$.
            The scheme 
            \begin{align}
                \X = \Proj{\OK[X,Y,Z]/(Y^3 - Z^3 - X^3 - 3XZ^2)}
            \end{align}
            has a non-reduced special fibre:
            On the affine chart $Z \neq 0$ we have $\X_s: 0 = y^3 - x^3 - 1 = (x - y - 1)^3$ over $\FF_3$.
            However, $\X$ is normal and thus a model of $X$ since
            \begin{align*}
                F(x,y) &= G(x,y)^r H_1(x,y) + 3 H_2(x,y) \\
                    &\coloneqq (x-y-1)^3 + 3(-x^2y + xy^2 + x^2 - 2xy + y^2 - y),
            \end{align*}
            i.e. $r=3$ but $e=1$ and $\bar{H}_2(x,y) \notin \bar{G}(x,y)\FF_3[x,y]$.
    \end{enumerate}
\end{example}

% Model with reduced special fibre: X: 0 = (y-1)(x^3+2) +3x

\section{From Models to Valuations}\label{sec:from-models-to-valuations}
In this section, we assume the valuation $v_K$ to be \textit{normalised}, i.e. $v_K(K^\times) = \ZZ$.

\begin{definition}
    Let $F_X$ be the function field of $X$ over $K$.
    A valuation $v: F_X \to \hat{\RR}$ is said to be \textit{of type \textrm{II}}, if
    \begin{enumerate}
        \item $v|_K = v_K$,
        \item $v$ is discrete, and
        \item the extension of residue fields $k(v)|k$ has transcendence degree $1$.
    \end{enumerate}
    A valuation of this type is also known as \textit{geometric valuation}.
    We write $V_\textrm{II}(F_X)$ for the set of all valuations of type \textrm{II} of $F_X$.
    The number
    \begin{align}
        m(v) \coloneqq [\Gamma_v : \Gamma_K]
    \end{align}
    is called the \textit{multiplicity of $v$}, where $\Gamma_v \coloneqq v(F_X^*)$ is the value group.
    We say that $v$ is \textit{weakly unramified}, if $m(v) = 1$.
\end{definition}

\begin{example}
    The \textit{Gauß valuation} on the polynomial ring $K[x]$ is defined as
    \begin{align}
        v_{\scriptscriptstyle G}: K[x] &\to \hat{\QQ}, \; v_{\scriptscriptstyle G}\left( \sum_i a_i x^{i} \right) = \min_i(v_K(a_i)).
    \end{align}
    Since $k(v_{\scriptscriptstyle G}) = k(\bar{x})$ (where $\bar{x}$ is the image of $x$ in $k(v_{\scriptscriptstyle G})$), the Gauß valuation
    extends naturally and uniquely to $K(x)$, meaning that we have $v_{\scriptscriptstyle G} \in V_\textrm{II}(K(x))$.
\end{example}

Let $\X$ be a model of $X$.
We wish to assign a valuation in $V_\textrm{II}(F_X)$ to each generic point of the special fibre $\X_s$.
Suppose that $Z$ is an irreducible component of $\X_s$ with generic point $\xi$.
Since $\xi$ has codimension $1$ in $\X$, the stalk $\OO_{\X,\xi}$ is a discrete valuation ring\@.
Thus, we can define a valuation on $F_X$,
\begin{align} \label{eq:ValuationOfComponent}
    v_Z: \mathrm{Frac}(\OO_{\X, \xi}) = F_X \to \hat{\QQ}.
\end{align}
Since $\X$ is a flat $\OK $-scheme $\OO_{\X, \xi}$ dominates $\OK$, and we may scale $v_Z$ such that the restriction of $v_Z$ to $K$ is $v_K$.
We set
\begin{align}
    V(\X) \coloneqq \lbrace v_Z: \; Z \subset \X_s \text{ is an irreducible component}\rbrace.
\end{align}

\begin{proposition}\label{prop:WellDefOfInducedValuation}
    The valuation defined in~\eqref{eq:ValuationOfComponent} is of type \textrm{II}.
\end{proposition}

\begin{proof}
    As mentioned above, we can assume $v_Z|_K = v_K$.
    Furthermore, we have a closed immersion $Z \hookrightarrow \X$
    by definition, inducing a surjection $\OO_{\X,\xi} \to \OO_{Z,\xi}$.
    Thus $\trdeg(k(v_Z)) = \codim(\OO_{Z,\xi}) = \codim(v_Z) = 1$, showing $v_Z \in V_\textrm{II}(F_X)$.
\end{proof}

\begin{definition}
    The \textit{multiplicity} of the irreducible component $Z$ of $\X_s$ containing the generic point $\xi$ is defined as
    \begin{align}
        m(Z) \coloneqq \length(\OO_{\X,\xi}).
    \end{align}
\end{definition}

\begin{proposition}\label{prop:multiplicities-coincide}
    It holds $m(v_Z) = m(Z)$.
\end{proposition}

\begin{proof}
    As the model $\X$ is normal, the stalk $\OO_{\X, \xi}$ is a discrete valuation ring\@.
    Let $\pi_\xi$ denote an uniformizing parameter, then
    \begin{align}
        \OO_{\X_s,\xi} = \OO_{\X,\xi}/(\pi_\xi).
    \end{align}
    Thus, the length of the stalk at $\xi$ coincides with the length of the chain of ideals generated
    by the uniformizing parameter of the quotient ring:
    Since $(\pi_\xi) = (\pi_{Z,\xi}^m)$, where $m = (\Gamma_{v_Z}:\Gamma_K)$, it holds $m(v_Z) = \length(\OO_{\X, \xi})$ as desired.
\end{proof}

The main goal of this chapter and the next is to establish the following theorem.

\begin{theorem}\label{thm:bijectionModelsAndComponents}
    The map
    \begin{align} \label{eq:mapFromModelToValuation}
        \X \mapsto V(\X).
    \end{align}
    defines a monotone\footnote{cf. Lemma \ref{lemma:induced-order-on-valuations}} bijection between
    \begin{itemize}
        \item the set of isomorphism classes of models of $X$, and
        \item the set of finite, nonempty subsets of $V_\textrm{II}(F_X)$.
    \end{itemize}
\end{theorem}

In the remainder of this section, we prove injectivity of this map.
We follow \cite[Chapter 3]{Rueth}

\begin{lemma}\label{lem:preimageOfGenericPoint}
    Let $\X$ be a model of $X$ and let $\X'$ be a model of another irreducible smooth projective curve $X'$.
    Let $\varphi: \X' \to \X$ be a morphism of $\OK $-schemes which induces a cover $X' \to X$ on the generic fibre.
    Let $\eta$ be a generic point of $\X_s$, then the preimage of $\eta$ under $\varphi$ is non-empty and consists of generic points of $\X_s'$.
    If $\varphi$ is birational, then the preimage consists of precisely one generic point.
\end{lemma}

\begin{proof}{\cite[Lemma 3.6]{Rueth}}
    Let $Z \subset \X_s$ be the irreducible component that contains $\eta$.
    Since $\varphi$ is surjective (cf.\ Remark~\ref{rem:modifications-are-surjective-and-partial-order-of-models}),
    there is a point $\eta'$ on an irreducible component $Z' \subset \X_s'$ which is mapped to $\eta$.
    Since $\varphi$ is dominant, the induced ring homomorphism $\OO_{\X,\eta} \to \OO_{\X',\eta'}$ is injective.
    This is an extension of discrete valuation rings and thus $\eta' \in Z'$ must be generic.
    In particular, if $\varphi$ is birational, then both rings are integrally closed in the same field of fractions which makes them equal.
    Therefore, $\eta'$ must be unique since $\X'$ is separated.
\end{proof}

\begin{lemma}[{\cite[Lemma 3.7]{Rueth}}]\label{lemma:induced-order-on-valuations}
    Let $\varphi: \X' \to \X$ be a modification of models of $X$.
    Then $V(\X) \subset V(\X')$.
    In particular, the partial order $\X \leq \X'$ of models on $X$ induces a partial order on the corresponding finite
    subsets of $V_{\mathrm{II}}(F_X)$ via inclusion maps.
\end{lemma}

\begin{proof}
    This is a consequence of Proposition~\ref{prop:WellDefOfInducedValuation} and Lemma~\ref{lem:preimageOfGenericPoint}.
\end{proof}

\begin{proposition}[{\cite[Lemma 3.8]{Rueth}}]
    Let $\X$ and $\X'$ be models of $X$ with $V(\X) \subset V(\X')$.
    Then there is a modification $\X' \to \X$.
\end{proposition}

\begin{proof}
    If $V(\X) \subset V(\X')$, then for every valuation $v_Z \in V(\X)$ corresponding to an irreducible component $Z\subset \X_s$ with generic point $\xi$,
    there is an irreducible component $Z' \subset \X'_s$ with generic point $\xi'$, such that $\OO_{\X,\xi} = \OO_{\X',\xi'} \subset F_X$.
    Together with the identification $\X_g = X = \X'_g$, this induces a birational map $\X' \dashrightarrow \X$ (cf. for example \cite[Lemma 29.50.2.]{Stacks}).
    By \cite[Theorem 8.3.20]{Liu}, this map is a morphism. Since it identifies $\X_g$ and $\X'_g$, it is a modification.
\end{proof}


\begin{comment}
    \begin{proof}[Sketch of the proof]
    Let
    \begin{align*}
        U= \lbrace Z_\eta \subset \X_s: v_{Z_\eta} \in V(\X)\rbrace
    \end{align*}
    and
    \begin{align*}
        U'= \lbrace Z_{\eta'} \subset \X_s: v_{Z_{\eta'}} \in V(\X')\rbrace
    \end{align*}
    be covers of $\X_s$ (resp. $\X_s'$) of irreducible components, where the irreducible component
    $Z_\eta$ has generic point $\eta$.
    The inclusion $V(\X) \subset V(\X')$ induces an injective map $\iota$ from the generic points of the irreducible components of $\X_s$ to those of $\X'_s$, which induces
    ring homomorphisms $\OO_{\X_s', \iota(\eta)} \to \OO_{\X_s, \eta}$.
    Thus, we get a birational map $\X' \dbkarow \X$ which is a morphism by~\cite[Theorem 8.3.20]{Liu}.
\end{proof}
\end{comment}


\begin{corollary}[{\cite[Proposition 3.9]{Rueth},\cite[Corollary 8.3.23]{Liu}}]
    Two models $\X$ and $\X'$ of $X$ are isomorphic as models of $X$ if and only if $V(\X) = V(\X')$.
    In particular, the map defined in~\eqref{eq:mapFromModelToValuation} is injective.
\end{corollary}

\begin{proof}
    Suppose that $V(\X) = V(\X')$.
    Then there exist morphisms $\X\to \X'$ and $\X' \to \X$.
    The composition of these morphisms belongs to the equivalence class of the identity morphism, when considered as a rational map.
\end{proof}

\begin{example}
    We continue with Example~\ref{ex:model-of-plane-curve}.
    Let $F\in \QQ_p[X,Y,Z]$ be irreducible, integral and primitive and such that
    \begin{align}
        \X = \Proj{\ZZ_p[X,Y,Z]/(F)}
    \end{align}
    is a model of the projective curve $X: 0=F(X,Y,Z)$ over $\QQ_p$.
    Suppose that $\bar{F} \in \FF_p[X,Y,Z]$ is irreducible, then $\X_s$ has precisely one irreducible component $Z$
    with generic point $\xi$ corresponding to the ideal $(p, \bar{F})$.
    In particular, it holds $V(\X) = \lbrace v_Z\rbrace$. 
    $p$ is a unifomising parameter of $\OO_{\X, \eta}$, thus $v_Z(p)=1$ and since $x \notin \m_\eta=(p)$, $v_Z(x)=0$.
    Hence, on the function field $F_X$, $v_Z$ is given by 
    \begin{align}
        v_Z(f) = \min_{i,j}v_p(a_{i,j}) \quad \text{for } f = \sum_{i,j}a_{i,j}x^i y^j.
    \end{align}
    In particular $v_Z|_{\QQ_p(x)} = v_{\scriptscriptstyle G}$.
\end{example}

\begin{comment}
    \begin{example}
    We continue with Example~\ref{ex:particular-model-of-plane-curve} of the model
    \begin{align*}
        \mathcal{X} = \Proj \ZZ_p[X,Y,Z]/(F)
    \end{align*}
    with $F = Y^2 Z - X^3 - pXZ^2 \in \QQ_p[X,Y,Z]$ of the elliptic curve $X: y^2 = x^3 +px$ defined over $\QQ_p$.
    We have already checked that it is a model.
    Since $\bar{F} \in \FF_p[X,Y,Z]$ is irreducible, $\X_s$ consists of one irreducible component $Z$ with generic point $\xi$
    corresponding to the ideal $(p, Y^2Z - X^3)$.
    In particular, $V(\XX)$
    $p$ is a uniformising parameter of $\OO_{\X, \eta}$, thus $v_Z(p) = 1$ and since $x\notin \m_\eta = (p)$, $v_Z(x)=0$.
    Hence, on the function field $F_X = \QQ_p(x)[y]/(y^2 - x^3 - px)$, $v_Z$ is given by 
    \begin{align}
        v_Z(f) = \min_{i,j}v_p(a_{i,j}) \quad \text{for } f = \sum_{i,j}a_{i,j}x^i y^j.
    \end{align}
    In particular $v_Z|_{\QQ_p(x)} = v_{\scriptscriptstyle G}$.
\end{example}
\end{comment}


